\chapter*{\MakeUppercase{Lampiran}}
\addcontentsline{toc}{chapter}{\MakeUppercase{Lampiran}}

Lampiran berisi materi pendukung yang digunakan dalam pelaksanaan atau penyusunan penelitian, tetapi tidak ditempatkan pada bagian utama naskah agar pembahasan tetap ringkas dan terfokus. Materi yang disajikan dalam lampiran harus tetap relevan dengan penelitian dan dapat membantu pembaca memahami, memverifikasi, atau mereproduksi hasil penelitian.

Dokumen yang dapat disertakan dalam lampiran antara lain kode sumber program, data mentah, tabel yang berukuran besar, hasil perhitungan rinci, instrumen penelitian, hasil pengujian tambahan, dokumentasi kegiatan, serta dokumen pendukung lainnya yang diperlukan. Materi yang disajikan dalam lampiran sebaiknya merupakan informasi yang penting untuk melengkapi penelitian, namun kurang sesuai apabila ditempatkan pada bab-bab utama Tugas Akhir karena dapat mengganggu alur pembahasan.
